% 04_belief_loop_d1.tex

\section{D1-Closed Belief Loop}
\label{sec:belief-loop}

The previous section established two ingredients. First, conditional on a fixed judgment-intensity state, the belief-dependent payoff is continuous in the relevant higher-order belief hierarchy. Second, the exposure term is pinned to the sender's true type through a primitive ranking, \(\Phi_H>\Phi_L\). It is not a signal-precision parameter and not an equilibrium posterior.

This section separates the main argument into two independent reversals. The first is an IC reversal: once \(q>q_{\mathrm{IC}}\), \(\theta_H\) exits \(s_A\) even if the receiver would still read \(s_A\) as high-type evidence. The second is a D1 reversal: once \(q>q_{\mathrm{D1}}\), the high type has a strictly smaller profitable-deviation response set than the low type, so an off-path \(s_A\) is no longer high-type evidence but low-type evidence. The final theorem only combines these two reversals. For \(q>\max\{q_{\mathrm{IC}},q_{\mathrm{D1}}\}\), safe pooling is the unique D1-closed pure assessment.

\paragraph{Payoff convention.}
Fix a low-scrutiny history and write
\[
q=\Pr(j_{t+1}=H\mid j_t=L)
\]
for the transition probability into high scrutiny. For type \(\theta\in\{\theta_L,\theta_H\}\), let
\[
B_\theta(\rho)
\]
denote the current net benefit from choosing \(s_A\) rather than \(s_P\) when the ambitious signal is interpreted with posterior \(\rho\). This term includes the current reputational return from the receiver's interpretation, net of current signal cost and low-scrutiny epistemic cost. It excludes future high-scrutiny exposure. Let \(\Phi_\theta>0\) denote the true-type-indexed adverse-repricing exposure cost, with the reversed ranking \(\Phi_H>\Phi_L\) supplied by Section~\ref{sec:psychological-exposure}.

Let \(e_t\ge 0\) denote exposed ambition at the beginning of period \(t\). Sending \(s_A\) adds one unit of new exposed ambition; choosing \(s_P\) adds no new unit. Existing exposure depreciates at rate \(\eta\in(0,1]\):
\[
e_{t+1}=(1-\eta)e_t+\mathbf{1}\{s_t=s_A\}.
\]
Let
\[
\Gamma(q,\eta)>0
\]
denote the discounted marginal continuation price of adding one unit of exposed ambition in the current low-scrutiny state. The one-period reduced form is the special case
\[
\Gamma(q,\eta)=\beta q\mu_H,
\]
where \(\beta\in(0,1)\) is the discount factor and \(\mu_H\) is the high-scrutiny intensity. In the stock version,
\[
\Gamma(q,\eta)
=
\mathbb E_q\left[
\sum_{\tau\ge 1}
\beta^\tau (1-\eta)^{\tau-1}
\mu_H\mathbf{1}\{j_{t+\tau}=H\}
\,\middle|\,j_t=L
\right].
\]
Low-scrutiny and current epistemic costs are already absorbed into \(B_\theta(\rho)\); \(\Gamma(q,\eta)\) prices only the marginal future adverse-repricing exposure created by entering high scrutiny. The maintained monotonicity condition is
\[
\Gamma(q,\eta)\text{ is continuous and strictly increasing in }q.
\]

The payoff margin used throughout this section is
\[
\Delta_\theta(\rho,q)
=
B_\theta(\rho)-\Gamma(q,\eta)\Phi_\theta.
\]
The fixed high-attribution surplus sometimes written \(B_H\) is \(B_H(1)\).

\subsection{IC Reversal}
\label{subsec:ic-reversal}
\label{subsec:anticipatory-incentive-failure}

At the low-scrutiny history, the sender chooses between the ambitious signal \(s_A\) and the pooling-safe signal \(s_P\). Sending \(s_A\) creates an exposure state that can be repriced if the next-period state is \(H\).

The high type's separating incentive constraint is
\[
\Delta_H(1,q)\ge 0.
\]
Thus define
\[
q_{\mathrm{IC}}
=
\inf\{q:\,B_H(1)-\Gamma(q,\eta)\Phi_H<0\}.
\]
In the one-period exposure specification, this collapses to
\[
q_{\mathrm{IC}}=\frac{B_H(1)}{\beta\mu_H\Phi_H},
\]
assuming \(B_H(1)>0\). For \(q>q_{\mathrm{IC}}\), the high type strictly prefers the pooling-safe signal \(s_P\) to remaining on the ambitious signal \(s_A\), even if \(s_A\) is still interpreted as high type.

\begin{proposition}[IC reversal]
\label{prop:ic-reversal}
\label{lem:anticipatory-incentive-failure}
\label{lem:general-price-ic-reversal}
Suppose the current state is \(L\), \(B_H(1)>0\), \(\Phi_H>0\), and \(\Gamma(q,\eta)\) is continuous and strictly increasing in \(q\). If
\[
q>q_{\mathrm{IC}},
\]
then
\[
\Delta_H(1,q)=B_H(1)-\Gamma(q,\eta)\Phi_H<0.
\]
Hence
\[
q>q_{\mathrm{IC}}
\quad\Rightarrow\quad
\theta_H \text{ exits } s_A
\text{ even under the high-type interpretation } \rho_A=1.
\]
In particular, the separating profile
\[
\theta_H\mapsto s_A,
\qquad
\theta_L\mapsto s_P
\]
is not an equilibrium.
\end{proposition}

\begin{proof}
Under the separating profile, \(s_A\) is interpreted as the high-type signal. Hence the high type receives the current net benefit \(B_H(1)\) from \(s_A\), but also internalizes the future exposure cost \(\Gamma(q,\eta)\Phi_H\). Its payoff gain from choosing \(s_A\) rather than \(s_P\) is therefore
\[
\Delta_H(1,q)=B_H(1)-\Gamma(q,\eta)\Phi_H.
\]
Since \(\Phi_H>0\) and \(\Gamma(q,\eta)\) is strictly increasing in \(q\), this expression is strictly decreasing in \(q\). By the definition of \(q_{\mathrm{IC}}\), every \(q>q_{\mathrm{IC}}\) lies in the region where \(B_H(1)-\Gamma(q,\eta)\Phi_H<0\). Thus the high type strictly prefers \(s_P\) to \(s_A\). The proposed separating profile violates the high type's incentive constraint and therefore cannot be an equilibrium.
\end{proof}

The point is not that the low type successfully mimics the high type. The separating profile collapses before that question arises. The IC reversal says that high-type meaning is no longer enough to keep the high type on the ambitious signal: the expected future cost of adverse repricing exceeds the current benefit of being recognized.

\subsection{D1 Reversal}
\label{subsec:d1-reversal}
\label{subsec:d1-off-path-reversal}

It remains to check whether pooling on \(s_P\) can be sustained. The possible obstruction is off-path resurrection: if all types pool on \(s_P\), an off-path deviation to \(s_A\) might again be interpreted as evidence of high type, restoring the value of the ambitious signal.

This is where the refinement matters. Consider the pooling candidate in which both types choose \(s_P\). The receiver observes an off-path signal \(s_A\). Let \(a\) denote a receiver response, and let \(\rho(a)\) be the posterior or interpretation induced by that response. D1 determines how the receiver interprets the off-path signal and therefore changes the current payoff term \(B_\theta(\rho(a))\). It does not relabel the deviator's true-type exposure technology \(\Phi_\theta\).

Define the profitable-deviation response set
\[
\mathcal D_\theta(q)=\{a:\Delta_\theta(\rho(a),q)>0\}.
\]

In the standard Spence ranking, the high type has a lower signaling cost, so the affine static case gives
\[
\mathcal D_L(q)\subset \mathcal D_H(q).
\]
D1 then assigns the off-path ambitious signal to the high type. The present model reverses this inclusion when future scrutiny is sufficiently likely. By Lemma~\ref{lem:bayesian-surprise-ordering}, high types have larger future exposure:
\[
\Phi_H>\Phi_L.
\]
Let
\[
M:=\sup_{a}\{B_H(\rho(a))-B_L(\rho(a))\},
\]
and assume \(M<\infty\). Define the generalized D1 threshold
\[
q_{\mathrm{D1}}
=
\inf\{q:\Gamma(q,\eta)(\Phi_H-\Phi_L)>M\}.
\]
For \(q>q_{\mathrm{D1}}\), every receiver response satisfies
\[
\Delta_H(\rho(a),q)<\Delta_L(\rho(a),q),
\]
because
\[
\Delta_H(\rho(a),q)-\Delta_L(\rho(a),q)
=
B_H(\rho(a))-B_L(\rho(a))
-\Gamma(q,\eta)(\Phi_H-\Phi_L).
\]

\begin{proposition}[D1 reversal]
\label{prop:d1-reversal}
\label{lem:d1-off-path-reversal}
Suppose \(\Phi_H>\Phi_L\), \(M<\infty\), \(q>q_{\mathrm{D1}}\), and the receiver-response space is rich enough that there exists a response \(a\) with
\[
\Delta_L(\rho(a),q)>0>\Delta_H(\rho(a),q).
\]
Then
\[
q>q_{\mathrm{D1}}
\quad\Rightarrow\quad
\mathcal D_H(q)\subsetneq \mathcal D_L(q).
\]
Therefore, under the D1 criterion, an off-path ambitious signal \(s_A\) cannot be assigned to the high type. It is assigned to the low type:
\[
\rho_A^{D1}=0.
\]
\end{proposition}

\begin{proof}
If \(q>q_{\mathrm{D1}}\), then for every receiver response \(a\),
\[
\Delta_H(\rho(a),q)<\Delta_L(\rho(a),q).
\]
Therefore, whenever the deviation is profitable for the high type, it is also profitable for the low type:
\[
\Delta_H(\rho(a),q)>0
\quad\Rightarrow\quad
\Delta_L(\rho(a),q)>0.
\]
Thus
\[
\mathcal D_H(q)\subseteq\mathcal D_L(q).
\]
The strict-response richness condition gives some \(a\) satisfying \(\Delta_L(\rho(a),q)>0>\Delta_H(\rho(a),q)\), so \(a\in\mathcal D_L(q)\setminus\mathcal D_H(q)\). Hence
\[
\mathcal D_H(q)\subsetneq\mathcal D_L(q).
\]
Under D1, if one type can profitably deviate under a strictly wider set of receiver responses, the deviation is attributed to that type rather than to the type whose profitable-deviation set is smaller. Hence the off-path signal \(s_A\) is attributed to \(\theta_L\), not to \(\theta_H\), and the off-path posterior is \(\rho_A^{D1}=0\).
\end{proof}

The affine Spence payoff is a special case, not the maintained payoff notation. If
\[
B_\theta(\rho(a))=R(a)-\bar k_\theta,
\]
then
\[
M=\bar k_L-\bar k_H.
\]
In the one-period reduced form \(\Gamma(q,\eta)=\beta q\mu_H\), and when \(\bar k_L>\bar k_H\), the generalized threshold becomes the old closed form
\[
q_{\mathrm{D1}}
=
\frac{\bar k_L-\bar k_H}{\beta\mu_H(\Phi_H-\Phi_L)}.
\]

The ambitious signal therefore changes its meaning. In the standard model, an off-path ambitious signal says: this must be a strong type, because only strong types can afford it. Here, when future scrutiny is sufficiently likely, the same signal says: this is unlikely to be the high type, because high types have the most to lose from carrying exposed ambition into future scrutiny. This is the second reversal, independent of the IC reversal above.

\subsection{Belief Hierarchies and Closed Assessments}
\label{subsec:belief-loop-hierarchy}

The fixed point is taken over a closed assessment, not over a strategy profile alone. At a low-scrutiny history \(h_t\), let
\[
\sigma_\theta(h_t)\in\{s_A,s_P\}
\]
denote the type-\(\theta\) sender's signal choice. Let
\[
\rho_A(h_t):=\Pr(\theta_H\mid s_A,h_t)
\]
denote the receiver's first-order posterior after observing the ambitious signal. If \(s_A\) is on path, \(\rho_A\) is pinned down by Bayes' rule. If \(s_A\) is off path, \(\rho_A\) is selected by the refinement.

The sender's payoff depends on what she expects the receiver to believe after \(s_A\). Write
\[
\widehat \rho_{\theta,A}(h_t)
\]
for the type-\(\theta\) sender's second-order belief about the receiver's posterior following \(s_A\). A closed assessment is a triple
\[
\mathcal A=(\sigma,\rho_A,\widehat \rho_A),
\]
where \(\widehat \rho_A=(\widehat\rho_{H,A},\widehat\rho_{L,A})\). Belief consistency requires
\[
\widehat\rho_{\theta,A}=\rho_A
\qquad\text{for each }\theta\in\{\theta_H,\theta_L\}.
\]
Thus the sender's second-order belief about the receiver's first-order posterior must agree with the posterior induced by the assessment itself.

Crucially, the fixed point is evaluated exactly at this second-order hierarchy layer. Continuity of expected utility with respect to the relevant belief hierarchy ensures the belief-dependent payoff is regular within each regime slice. Here, the net benefit \(B_\theta(\widehat{\rho}_{\theta,A})\) depends continuously on the expected first-order posterior. While the D1 refinement injects a discontinuous jump into the receiver's belief map \(\rho_A(\sigma)\) when the signal becomes off path, the fixed point is resolved on the composition of these expectations, so the belief loop closes at \(\widehat\rho_{\theta,A}=\rho_A\).

This is the level at which the fixed point is closed. The exposure term \(\Phi_\theta\) remains true-type indexed. It is not changed by the receiver's posterior. The receiver's posterior changes the current reputational benefit of the signal; it does not relabel the deviator's adverse-repricing exposure.

Using the payoff convention above, the proof below only needs the current-benefit values at three posteriors: high-type attribution, prior attribution, and low-type attribution.
Let
\[
\pi_0:=\Pr(\theta_H)\in(0,1)
\]
denote the prior probability of the high type. Assume
\[
B_H(1)>0,
\qquad
B_H(\pi_0)\le B_H(1),
\qquad
B_H(0)\le 0.
\]
The middle inequality only rules out the possibility that all-ambitious pooling creates a larger high-type current benefit than separation. It is a local prior condition, not a monotonicity assumption over all posteriors. The last inequality says that an ambitious signal interpreted as low-type evidence carries no high-type premium. For the low type, assume
\[
B_L(0)<0.
\]
Thus if \(s_A\) is interpreted as low-type evidence, the low type does not receive enough current benefit to justify choosing it.

\subsection{The Belief-Loop Operator}
\label{subsec:belief-loop-operator}

For \(q>q_{\mathrm{D1}}\), the D1-selected off-path posterior after \(s_A\) is
\[
\rho_A^{D1}=0
\]
under the strict-response richness condition of Proposition~\ref{prop:d1-reversal}. This D1 conclusion is an off-path counterfactual rule. It is obtained by comparing the response sets \(\mathcal D_H(q)\) and \(\mathcal D_L(q)\) defined from \(\Delta_\theta(\rho(a),q)\); it does not presuppose that either type actually deviates.

Let
\[
z\in\{0,1\}
\]
denote whether the high type is expected to occupy the ambitious signal:
\[
z=1 \quad\text{means } \theta_H \text{ chooses }s_A,
\qquad
z=0 \quad\text{means } \theta_H \text{ withdraws to }s_P.
\]
For \(q>q_{\mathrm{D1}}\), the receiver's belief about \(s_A\) is
\[
\rho_A(z)=
\begin{cases}
1, & \text{if }z=1\text{ and }s_A\text{ is the separating high-type signal},\\
0, & \text{if }z=0\text{ and }s_A\text{ is off path, selected by D1}.
\end{cases}
\]
The first line is Bayes' rule on the separating path. The second line is the D1 off-path counterfactual map.

Given this belief map, the high type's participation response is
\[
T_q(z)
=
\mathbf{1}\left\{
\Delta_H(\rho_A(z),q)\ge 0
\right\}.
\]
A high-type belief-loop fixed point is a value \(z\in\{0,1\}\) such that
\[
z=T_q(z).
\]
This is the self-referential equation. High-type participation determines whether \(s_A\) is on path. The path status of \(s_A\) determines the receiver's posterior. The receiver's posterior determines the current rent of \(s_A\). That rent determines the sender's second-order belief and high-type participation.

Define
\[
q^*=\max\{q_{\mathrm{IC}},q_{\mathrm{D1}}\}.
\]
The closed-loop candidate that will survive is the safe-pooling configuration
\[
\mathcal C_0
=
\left(
\sigma_H=s_P,\,
\sigma_L=s_P;\,
s_A\text{ off path};\,
\rho_A^{D1}=0,\,
\widehat\rho_{H,A}=\widehat\rho_{L,A}=0
\right).
\]
It closes the loop because both types choose \(s_P\), D1 evaluates the off-path \(s_A\), D1 selects \(\rho_A^{D1}=0\), and sender second-order beliefs agree with that posterior.

\subsection{Safe Pooling as the Unique D1-Closed Pure Assessment}
\label{subsec:top-down-pooling-closed-fixed-point}

\begin{theorem}[Safe pooling as the unique D1-closed pure assessment]
\label{thm:top-down-pooling}
Suppose:
\begin{enumerate}
\item \(\pi_0=\Pr(\theta_H)\in(0,1)\), with \(B_H(1)>0\), \(B_H(\pi_0)\le B_H(1)\), and \(B_H(0)\le 0\);
\item \(B_L(0)<0\);
\item \(\Phi_H>\Phi_L\), with one primitive foundation given in Section~\ref{sec:psychological-exposure};
\item \(\Gamma(q,\eta)\) is continuous and strictly increasing in \(q\), with \(\eta\in(0,1]\);
\item \(q>q^*=\max\{q_{\mathrm{IC}},q_{\mathrm{D1}}\}\);
\item off-path beliefs satisfy D1, with the strict-response richness condition of Proposition~\ref{prop:d1-reversal};
\item the analysis is restricted to the binary-signal pure-strategy class \(\{s_A,s_P\}\).
\end{enumerate}
Then
\[
q>\max\{q_{\mathrm{IC}},q_{\mathrm{D1}}\}
\quad\Rightarrow\quad
\mathcal C_0 \text{ is the unique D1-closed pure assessment.}
\]
Equivalently, no D1-consistent closed assessment in this class has stable top-end occupancy of \(s_A\), and the unique D1-consistent pure closed-loop assessment is \(\mathcal C_0\):
\[
\theta_H\mapsto s_P,
\qquad
\theta_L\mapsto s_P.
\]
Equivalently, the unique high-type occupancy fixed point is
\[
z^*=0.
\]
\end{theorem}

\begin{proof}
We evaluate all pure closed-loop configurations, using the two reversals separately.

First, the IC reversal rules out top-end occupancy of \(s_A\). If only the high type occupies \(s_A\), then \(s_A\) is the separating high-type signal, so Bayes' rule gives
\[
\rho_A(1)=1.
\]
The high type's payoff gain from choosing \(s_A\) rather than \(s_P\) is
\[
\Delta_H(1,q)=B_H(1)-\Gamma(q,\eta)\Phi_H.
\]
Since \(q>q_{\mathrm{IC}}\), Proposition~\ref{prop:ic-reversal} makes this expression strictly negative. Hence \(T_q(1)=0\), so \(z=1\) is not a fixed point.

If both types choose \(s_A\), then \(s_A\) is pooling rather than separating. Bayes' rule gives \(\rho_A=\pi_0\). By assumption,
\[
B_H(\pi_0)\le B_H(1).
\]
Because \(q>q_{\mathrm{IC}}\),
\[
B_H(1)-\Gamma(q,\eta)\Phi_H<0.
\]
Therefore
\[
B_H(\pi_0)-\Gamma(q,\eta)\Phi_H<0.
\]
The high type would prefer \(s_P\), so pooling on \(s_A\) is not a fixed point either. Thus every configuration with top-end occupancy of \(s_A\) is infeasible.

Second, no low-only occupancy configuration is feasible. If only the low type chooses \(s_A\), then \(s_A\) is on path as low-type evidence, so \(\rho_A=0\). The low type's gain is
\[
B_L(0)-\Gamma(q,\eta)\Phi_L<0,
\]
because \(B_L(0)<0\) and \(\Gamma(q,\eta)\Phi_L>0\). Hence the low type does not wish to occupy \(s_A\) under low-type attribution.

Third, evaluate \(\mathcal C_0\). Under \(\mathcal C_0\), \(s_A\) is off path. Since \(q>q_{\mathrm{D1}}\), Proposition~\ref{prop:d1-reversal} implies
\[
\mathcal D_H(q)\subsetneq \mathcal D_L(q),
\]
so the off-path ambitious signal is assigned to the low type and
\[
\rho_A^{D1}=0.
\]
This is the load-bearing use of D1 in the theorem: the off-path ambitious deviation is evaluated under low-type attribution rather than under an unrestricted resurrection belief.
Belief consistency then sets
\[
\widehat\rho_{H,A}=\widehat\rho_{L,A}=\rho_A^{D1}=0.
\]
The high-type deviation margin is
\[
\Delta_H(0,q)=B_H(0)-\Gamma(q,\eta)\Phi_H<0,
\]
and the low-type deviation margin is
\[
\Delta_L(0,q)=B_L(0)-\Gamma(q,\eta)\Phi_L<0.
\]
Thus neither type has a profitable deviation to \(s_A\). The safe-pooling configuration is internally consistent. Since all other pure configurations have been ruled out, \(\mathcal C_0\) is the unique D1-consistent closed assessment in the binary pure-strategy class.
\end{proof}

The theorem should be read as a state claim:
\[
\mathcal C_0
\Longleftrightarrow
\left[
s_A\text{ off path}
\right]
\wedge
\left[
\rho_A^{D1}=0
\right]
\wedge
\left[
\widehat\rho_{H,A}=\widehat\rho_{L,A}=0
\right]
\wedge
\left[
\Delta_H(0,q)<0
\right]
\wedge
\left[
\Delta_L(0,q)<0
\right].
\]
The D1 component is not caused by an observed low-type deviation. It is the refinement's counterfactual response to an off-path signal, determined by the D1 reversal \(\mathcal D_H(q)\subsetneq\mathcal D_L(q)\). The participation inequalities then show that the counterfactual off-path interpretation is consistent with nobody actually sending \(s_A\).

This separation removes the circular timing story. The IC reversal is the high-type inequality \(\Delta_H(1,q)<0\) once \(q>q_{\mathrm{IC}}\): \(\theta_H\) exits \(s_A\) even under high-type interpretation. The D1 reversal is the response-set inclusion \(\mathcal D_H(q)\subsetneq\mathcal D_L(q)\) once \(q>q_{\mathrm{D1}}\): off-path \(s_A\) becomes low-type evidence rather than high-type evidence. The theorem combines them only at \(q>\max\{q_{\mathrm{IC}},q_{\mathrm{D1}}\}\), where safe pooling is the unique D1-closed pure assessment.

In the coordination and self-fulfilling-belief literature, forward-induction refinement characteristically selects \emph{toward} efficiency: it eliminates the off-path beliefs that sustain coordination failure, uniquely selecting the Pareto-dominant equilibrium \citep[e.g., the refinement resolutions of][]{diamond1983bank}. This standard selection presumes that the party able to orchestrate a deviation is the one seeking the efficient outcome, so forward induction attributes the deviation's intent accordingly. We exhibit the opposite operation. When malicious scrutiny reverses the cost ranking so that the wider profitable-deviation set belongs to the \emph{low} type, the same D1 logic confirms rather than dissolves the absence of stable top-end occupancy: the refinement selects safe pooling. This resulting top-down pooling is an \emph{informational} disaster, destroying the market's ability to screen types. Yet it is also privately rational risk avoidance for high types facing excessive epistemic exposure. The refinement does not rescue the informational efficiency of the separating equilibrium; instead, it enforces the nonexistence of an occupied ambitious-top fixed point.

The theorem is deliberately pure-strategy. It does not use an interior mixing probability for the high type. In a binary model with no low-type background arrival at \(s_A\), any strictly positive high-type mass on \(s_A\) would keep the posterior after \(s_A\) degenerate in the high type. Interior mixing requires an additional object, such as low-type mimetic entry, noisy audience inference, trembles, or a continuous type distribution. Those objects reappear as the polarization and population-dynamics implications below. The present section proves only the sharper fixed-point claim: above \(q^*\), the only D1-consistent pure belief loop is \(\mathcal C_0\), so the ambitious signal has no stable top-end occupancy fixed point.
